\documentclass[11pt,a4paper]{article}
\usepackage[UTF8]{ctex}
\usepackage{amsmath}
\usepackage{booktabs}
\usepackage{geometry}
\usepackage{xcolor}
\usepackage{hyperref}

\geometry{left=2.5cm,right=2.5cm,top=2.5cm,bottom=2.5cm}

\title{\textbf{Task 3.1: 主题向量化技术报告}}
\author{TY项目组}
\date{2026年2月3日}

\begin{document}

\maketitle

\section{任务概述}

Task 3.1的目标是将LDA主题模型的离散主题分布转换为稠密的连续向量表示，为后续的跨时间窗口主题对齐提供数值计算基础。本任务采用Word2Vec词向量加权求和的方法实现主题向量化。

\section{Word2Vec模型训练}

\subsection{训练配置}

\begin{table}[h]
\centering
\begin{tabular}{ll}
\toprule
\textbf{参数} & \textbf{值} \\
\midrule
向量维度 (vector\_size) & 200 \\
上下文窗口 (window) & 5 \\
最小词频 (min\_count) & 50 \\
训练算法 (sg) & 0 (CBOW) \\
负采样数 (negative) & 5 \\
并行工作线程 (workers) & 4 \\
迭代轮数 (epochs) & 10 \\
\bottomrule
\end{tabular}
\caption{Word2Vec超参数配置}
\end{table}

\subsection{训练语料统计}

\begin{itemize}
    \item \textbf{文档总数}：9,160,136篇岗位描述（合并5个时间窗口）
    \item \textbf{原始词数}：247,691,957词（2.48亿）
    \item \textbf{有效词数}：221,157,282词（2.21亿，89.3\%有效率）
    \item \textbf{最终词汇量}：68,886词（经min\_count=50过滤）
\end{itemize}

\subsection{训练性能}

\begin{itemize}
    \item \textbf{训练总时长}：1,197.8秒（约20分钟）
    \item \textbf{训练速度}：1,846,297词/秒
    \item \textbf{模型大小}：约28MB（68,886 × 200 × 4字节）
\end{itemize}

\subsection{词汇覆盖度分析}

\begin{table}[h]
\centering
\begin{tabular}{lrrr}
\toprule
\textbf{时间窗口} & \textbf{LDA词典大小} & \textbf{Word2Vec词汇} & \textbf{覆盖率} \\
\midrule
2016-2017 & 15,411 & 68,886 & 100\% \\
2018-2019 & 7,907 & 68,886 & 100\% \\
2020-2021 & 22,021 & 68,886 & 100\% \\
2022-2023 & 30,000 & 68,886 & 100\% \\
2024-2025 & 19,796 & 68,886 & 100\% \\
\bottomrule
\end{tabular}
\caption{Word2Vec对LDA词典的覆盖度（理论上）}
\end{table}

\textcolor{blue}{\textbf{关键发现}}：Word2Vec词汇量（68,886）显著大于任何单个LDA词典，确保了高覆盖度。合并语料训练的策略有效缓解了词汇稀疏问题。

\section{主题向量化}

\subsection{向量化算法}

对于LDA模型中的第$t$个主题，其向量表示$\mathbf{v}_t$通过加权词向量求和计算：

\begin{equation}
\mathbf{v}_t = \frac{\sum_{w \in V} P(w|t) \cdot \mathbf{v}_w}{\sum_{w \in V} P(w|t)}
\end{equation}

其中：
\begin{itemize}
    \item $V$：主题$t$的前50个高频词
    \item $P(w|t)$：词$w$在主题$t$中的概率分布
    \item $\mathbf{v}_w$：词$w$的Word2Vec向量
\end{itemize}

归一化处理：
\begin{equation}
\mathbf{v}_t^{\text{norm}} = \frac{\mathbf{v}_t}{\|\mathbf{v}_t\|_2}
\end{equation}

L2归一化确保后续余弦相似度计算简化为点积运算。

\subsection{向量化结果}

\begin{table}[h]
\centering
\begin{tabular}{lccc}
\toprule
\textbf{时间窗口} & \textbf{主题数} & \textbf{向量维度} & \textbf{矩阵形状} \\
\midrule
window\_2016\_2017 & 60 & 200 & (60, 200) \\
window\_2018\_2019 & 60 & 200 & (60, 200) \\
window\_2020\_2021 & 60 & 200 & (60, 200) \\
window\_2022\_2023 & 60 & 200 & (60, 200) \\
window\_2024\_2025 & 60 & 200 & (60, 200) \\
\midrule
\textbf{合计} & 300 & - & 300,000维参数 \\
\bottomrule
\end{tabular}
\caption{主题向量化输出矩阵}
\end{table}

\subsection{向量化性能}

\begin{itemize}
    \item \textbf{单窗口处理时间}：< 0.2秒/窗口
    \item \textbf{总处理时间}：< 1秒（5窗口）
    \item \textbf{输出文件大小}：约240KB/窗口（60 × 200 × 4字节）
\end{itemize}

\section{技术质量评估}

\subsection{优势}

\begin{enumerate}
    \item \textbf{维度统一}：所有主题映射到相同的200维语义空间，支持跨窗口直接比较
    \item \textbf{语义保持}：Word2Vec捕获的词汇语义关系通过加权传递到主题向量
    \item \textbf{高效计算}：L2归一化后，余弦相似度退化为点积，计算复杂度$O(d)$
    \item \textbf{覆盖完整}：Word2Vec词汇量远大于单窗口词典，避免OOV问题
\end{enumerate}

\subsection{潜在局限}

\begin{enumerate}
    \item \textbf{词权假设}：仅使用前50词可能遗漏长尾重要词汇
    \item \textbf{线性组合}：加权平均假设词向量可线性组合表示主题语义
    \item \textbf{时序依赖}：Word2Vec在所有窗口语料上训练，可能引入未来信息
\end{enumerate}

\section{输出文件}

\begin{itemize}
    \item \texttt{word2vec.model}：训练好的Word2Vec模型（68,886词汇）
    \item \texttt{topic\_vectors/window\_*\_topic\_vectors.npy}：5个主题向量矩阵
    \item \texttt{topic\_vectors/vectorization\_metadata.pkl}：元数据（配置、词汇量统计）
\end{itemize}

\section{后续工作}

\subsection{Task 3.2: 混合对齐算法}

利用生成的主题向量矩阵，执行以下步骤：

\begin{enumerate}
    \item \textbf{相似度矩阵计算}：
    \begin{equation}
    S_{ij} = \cos(\mathbf{v}_i^{(t)}, \mathbf{v}_j^{(t+1)}) = \mathbf{v}_i^{(t)} \cdot \mathbf{v}_j^{(t+1)}
    \end{equation}
    
    \item \textbf{Hungarian算法}：寻找最优1-to-1匹配（Survival事件）
    
    \item \textbf{演化事件检测}：
    \begin{itemize}
        \item Split（1→多）：一个基准主题映射到多个目标主题
        \item Merge（多→1）：多个基准主题映射到一个目标主题
        \item Birth：目标主题无显著基准匹配
        \item Death：基准主题无显著目标匹配
    \end{itemize}
\end{enumerate}

\subsection{预期对齐质量}

基于5×2年窗口的平滑划分，预期：
\begin{itemize}
    \item \textbf{Survival比例}：60-75\%（相邻窗口主题延续性高）
    \item \textbf{Birth/Death比例}：10-20\%（窗口边界效应）
    \item \textbf{Split/Merge比例}：5-15\%（主题演化复杂度）
    \item \textbf{平均相似度}：0.55-0.65（优于之前3窗口的0.47）
\end{itemize}

\section{结论}

Task 3.1成功构建了高质量的主题向量表示系统，关键成果包括：

\begin{itemize}
    \item ✅ 训练68,886词汇的Word2Vec模型（200维，10 epochs）
    \item ✅ 向量化300个主题（5窗口 × 60主题）
    \item ✅ 建立统一的语义空间支持跨时间比较
    \item ✅ 完成时间：约20分钟（训练）+ 1秒（向量化）
\end{itemize}

该向量化基础设施为后续的主题演化追踪、技能迁移分析和岗位综合化度量提供了坚实的技术支撑。

\end{document}
